\input{preamble}

\title{Section 6.3 --- Circles --- Answer Key}

\newcommand{\crans}[2]{%
Center: $(#1)$; Radius: $#2$%
}

\newcommand{\ceq}[3]{%
$\left(x #1\right)^2+\left(y #2\right)^2 = #3$%
}

\begin{document}
\maketitle

\section*{Basic Equation of a Circle}
On a separate sheet of paper, sketch each circles described below.
\begin{smenum}
\mitemxx
{\begin{GridFilled}{6}{2}{3}{5}
\filldraw(-2,1) circle (\dotsize);
\draw[thick](-2,1) circle (4);
\end{GridFilled}}
{\begin{GridFilled}{0}{3}{0}{4}
\filldraw(2,3) circle (\dotsize);
\draw[thick](2,3) circle (1);
\end{GridFilled}}\end{smenum}
Give an equation of each circles described below.
\begin{smenum}
\mitemxx
{$(x-3)^2+(y-5)^2=4$}
{$(x+5)^2+(y+2)^2=36$}
\end{smenum}
Give an equation of the circle shown in each graph below.
\begin{smenum}
\mitemxx
{$(x-5)^2+(y+4)^2=9$}
{$(x+3)^2+(y-2)^2=25$}
\end{smenum}
Give the center and radius of the circle described by each equation below.
\begin{smenum}
\mitemxx
{\crans{5,-2}{7}}
{\crans{-6,-1}{2\sqrt{10}}}
\end{smenum}
On a separate sheet of graph paper, sketch the graph of each equation below.
\begin{smenum}
\mitemxx
{\begin{GridFilled}{7}{3}{2}{8}
\filldraw(-2,3) circle (\dotsize);
\draw[thick](-2,3) circle (5);
\end{GridFilled}}
{\begin{GridFilled}{6}{0}{4}{0}
\filldraw(-4,-2) circle (\dotsize);
\draw[thick](-4,-2) circle ({sqrt(3)});
\end{GridFilled}
Note: do your best to estimate the radius $\sqrt{3}\approx 1.732$}
\end{smenum}

\section*{Writing Equations of Circles}
Give the equation of each circle described below.
\begin{smenum}
\mitemxx
{\ceq{-4}{-1}{41}}
{\ceq{-2}{+1}{50}}
\mitemxx
{\ceq{+7}{-5}{13}}
{\ceq{+4}{+2}{20}}
\mitemxx
{\ceq{+2}{-4}{49}}
{\ceq{-3}{+5}{9}}
\mitemxx
{\ceq{+5}{-1}{1}}
{\ceq{+4}{+3}{34}}
\end{smenum}
Give an equation of the circles graphed below.
\begin{smenum}
\mitemxx
{\ceq{-2}{+1}{5}}
{\ceq{-2}{-2}{2}}
\end{smenum}

\section*{Working with Standard Form}
Give an equation of each circle below in standard form.
\begin{smenum}
\mitemxx
{$x^2+y^2-8x+4y+11=0$}
{$x^2+y^2+4x+10y+19=0$}
\mitemxx
{$x^2+y^2-3x-2y-13=0$}
{$x^2+y^2-4x-2y-4=0$}
\end{smenum}
Rewrite each equation below in center-radius form. If the equation describes a true circle, give its center and radius. If it doesn't, state whether it is the equation of a degenerate circle (single point), or if the equation has no solution.
\begin{smenum}
\mitemxx
{\crans{2,3}{2}}
{\crans{1,-3}{5}}
\mitemxx
{\crans{1,-2}{\frac{3}{2}}}
{\crans{0,-5}{4}}
\mitemxx
{Degenerate circle (single point $(3,5)$)}
{\crans{-2,4}{3}}
\mitemxx
{No solution}
{Degenerate circle (single point $(1,-5)$)}
\end{smenum}
Rewrite each equation below in center-radius form. Then, graph the circle on a separate sheet of graph paper.
\begin{smenum}
\mitemxx
{$x^2+(y-2)^2=4^2$\\
\begin{GridFilled}{4}{4}{2}{6}
\filldraw(0,2) circle (\dotsize);
\draw[thick](0,2) circle (4);
\end{GridFilled}}
{\ceq{+5}{+1}{3^2}\\
\begin{GridFilled}{8}{0}{4}{2}
\filldraw(-5,-1) circle (\dotsize);
\draw[thick](-5,-1) circle (3);
\end{GridFilled}}
\mitemx
{\ceq{-1}{+\frac{7}{2}}{3^2}\\
\begin{GridFilled}{2}{4}{7}{0}
\filldraw(1,-3.5) circle (\dotsize);
\draw[thick](1,-3.5) circle (3);
\end{GridFilled}}
\end{smenum}

\end{document}